% generated by R/04_tables.R from table1_group_models.md -- do not edit
% needs \usepackage{booktabs,adjustbox} in the preamble
\begin{table}[htbp]
\centering
\caption{Do the strain groupings explain variation in biofilm between strains?}%
\label{tab:biofilmmodels}
\small
\begin{adjustbox}{max width=\linewidth}
\begin{tabular}{lllrrr}
\toprule
\textbf{Strains} & \textbf{Model} & \textbf{Grouping} & \textbf{Groups} & \textbf{w} & \textbf{P} \\
\midrule
all (11) & 1 & global mean & 1 & 0.60 & -- \\
all (11) & 2 & ancestor vs evolved & 2 & 0.20 & 0.928 \\
all (11) & 3 & mutation vs none & 2 & 0.20 & 0.829 \\
evolved (10) & 1 & global mean & 1 & 0.05 & -- \\
evolved (10) & 4 & spore vs total & 2 & 0.02 & 0.773 \\
evolved (10) & 5a & sinR vs ywcC vs none & 3 & 0.93 & 0.028 \\
sinR/ywcC (8) & 1 & global mean & 1 & 0.01 & -- \\
sinR/ywcC (8) & 5 & sinR vs ywcC & 2 & 0.99 & 0.004 \\
\bottomrule
\end{tabular}
\end{adjustbox}
\par\smallskip
{\footnotesize\raggedright Linear mixed models of log blank-corrected OD550, one per grouping, with strain as a random effect. \textit{w} is the Akaike weight: the relative support for each grouping among those fit to the same strains, summing to one within each set. \textit{P} is a Kenward--Roger F-test against the global-mean model on the same strains. Sets of strains are not comparable with one another. Model 3 (mutation vs. no mutation) on all eleven strains compares the reference well, m23 and m26 --- none of which carries a mutation --- with the eight \textit{sinR} and \textit{ywcC} clones; among the ten evolved clones it is the same partition as model 4. Models 2 and 3 on all eleven strains include the reference well, recorded as \textit{B. subtilis} 168 $\Delta$6 and probably not this experiment's ancestor; models 4, 5 and 5a do not use it.\par}
\end{table}
